%------------------------------------------------------------------------------%
\documentclass[fleqn,utf8,aspectratio=169,12pt]{beamer}

\usepackage{integracao_e_series}

\title{Séries Numéricas -- Teste da Raiz}

%------------------------------------------------------------------------------%
\begin{document}

\addtitle

%------------------------------------------------------------------------------%
%------------------------------------------------------------------------------%
\section{Teste da Raiz}

%------------------------------------------------------------------------------%
\begin{frame}{Teste da Raiz}
	Seja $\;\sum a_n\;$ tal que $\;a_n\geq 0\;$ para $\;n\geq N\;$ \pause e
	\[
    \rho = \lim_{n\to \infty} \sqrt[n]{a_n}
  \]

  \pause
  Se $\rho<1$                    \tabto{37mm} a série \emph{converge}

  \pause
  Se $\rho>1$ ou $\rho\to\infty$ \tabto{37mm} a série \emph{diverge}

  \pause
  Se $\rho=1$                    \tabto{37mm} o teste é \emph{inconclusivo}
\end{frame}

%------------------------------------------------------------------------------%
\begin{frame}{Limite Importante}
	\begin{align*}
		\onslide<+->{\lim_{n\to\infty} \sqrt[n]{n}}
		\onslide<+->{& = \lim n^{\sfrac{1}{n}} \\[2mm]}
		\onslide<+->{& = \lim \exp \big( \ln n^{\sfrac{1}{n}} \,\big) \\[2mm]}
		\onslide<+->{& = \lim \exp \left( \frac{\ln n}{n} \right)}
	\end{align*}

	\[
		\onslide<+->{\lim_{x\to\infty} \frac{\ln x}{x}}
		\onslide<+->{\;=\; \lim \frac{\,\sfrac{1}{x}\,}{1}}
		\onslide<+->{\;=\; \lim \frac{1}{x}}
		\onslide<+->{\;=\; 0}
	\]

	\[
		\onslide<+->{\lim_{n\to\infty} \sqrt[n]{n}}
		\onslide<+->{\;=\; \exp \left( \lim \frac{\ln n}{n} \right)}
		\onslide<+->{\;=\; e^0}
		\onslide<+->{\;=\; 1}
	\]
\end{frame}

%------------------------------------------------------------------------------%
%------------------------------------------------------------------------------%
\section{Exemplos do Teste da Raiz}

%------------------------------------------------------------------------------%
\addtocounter{exemplo}{1}
\begin{frame}{Exemplo \theexemplo}
  Use o teste da raiz para analisar a convergência da série
  \[
  \sum_{n=1}^\infty \frac{4^n}{\,(3n)^n\,}
  \]
\end{frame}

%------------------------------------------------------------------------------%
\begin{frame}{Exemplo \theexemplo}
  \[
  a_n
  = \frac{4^n}{\,(3n)^n\,}
  \pause
  = \left(\frac{4}{\,3n\,}\right)^{\!n}
  \]

  \pause
  \[
  \sqrt[n]{a_n\,}
  \pause
  = \sqrt[n]{\left(\frac{4}{\,3n\,}\right)^{\!n}}
  \pause
  = \frac{4}{3n}
  \pause
  \to 0
  \pause
  < 1
  \]

  A série converge
\end{frame}
%------------------------------------------------------------------------------%

%------------------------------------------------------------------------------%
\addtocounter{exemplo}{1}
\begin{frame}{Exemplo \theexemplo}
	\begin{columns}
		\column{0.48\textwidth}
		\[
		a_n = \begin{cases}
			\displaystyle \;\dfrac{n}{\,2^n\,} & n \text{ impar} \\[5mm]
			\displaystyle \;\dfrac{1}{\,2^n\,} & n \text{ par}
		\end{cases}
		\]

		\pause
		\vspace{0.5\baselineskip}
		\[
		\sqrt[n]{a_n} = \begin{cases}
			\displaystyle \;\dfrac{\,\sqrt[n]{n}\,}{2} & n \text{ impar} \\[5mm]
			\displaystyle \;\;\,\dfrac{1}{\,2\,}       & n \text{ par}
		\end{cases}
		\]

		\column{0.48\textwidth}
		\pause
		\[
		\frac{1}{\,2\,} \;\leq\; \sqrt[n]{a_n} \;\leq\; \frac{\,\sqrt[n]{n}\,}{2}
		\]

		\pause
		Sabemos que \(\; \lim\sqrt[n]{n} = 1 \)

		\pause
		\vspace{.5\baselineskip}
		Pelo teorema do confronto
		\[\lim\sqrt[n]{a_n} = \frac{1}{2} < 1\]

		\pause
		Portanto a série é convergente

	\end{columns}
\end{frame}

%------------------------------------------------------------------------------%
\addtocounter{exemplo}{1}
\begin{frame}{Exemplo \theexemplo}
  Use o teste da raiz para analisar a convergência da série
  \[
    \sum_{n=1}^\infty\frac{1}{n^{n+1}}
  \]
\end{frame}

%------------------------------------------------------------------------------%
\begin{frame}{Exemplo \theexemplo}
  Fazendo a mudança de índice \(k=n+1\)
  \[
    \sum_{n=1}^\infty\frac{1}{\, n^{n+1}\,} =
    \sum_{k=2}^\infty\frac{1}{\,(k-1)^{k}\,}
  \]

  \pause
  \[
    a_k = \frac{1}{\,(k-1)^{k}\,}
  \]

  \pause
  \[
    \sqrt[k]{a_k}
    = \frac{1}{k-1}
    \pause
    \to 0
    \pause
    < 1
  \]

  A série converge
\end{frame}

%------------------------------------------------------------------------------%
\addtocounter{exemplo}{1}
\begin{frame}{Exemplo \theexemplo}
  Use o teste da raiz para analisar a convergência da série
  \[
    \sum_{n=1}^\infty \left[\ln\left(e^2+\frac{1}{n}\right)\right]^{n+1}
  \]
\end{frame}

%------------------------------------------------------------------------------%
\begin{frame}{Exemplo \theexemplo}
  Fazendo a mudança de índice \(k=n+1\)
  \[
    \sum_{n=1}^\infty \left[\ln\left(e^2+\frac{1}{n}  \right)\right]^{n+1} \;=\;
    \sum_{k=2}^\infty \left[\ln\left(e^2+\frac{1}{k-1}\right)\right]^{k}
  \]

  \pause
  \[
    a_k
    = \left[\ln\left( e^2 + \frac{1}{k-1} \right)\right]^k
    \pause
    \qquad
    \sqrt[k]{a_k}
    = \ln\left(e^2+\frac{1}{k-1}\right)
  \]

  \pause
  Aplicando o Teste da Raiz
  \[
    \rho
    = \lim \sqrt[k]{a_k}
    \pause
    = \lim \ln\left(e^2+\frac{1}{k-1}\right)
    \pause
    = \ln \left(e^2+ \lim\frac{1}{k-1}\right)
    \pause
    = \ln \left(e^2\right)
    \pause
    = 2
    \pause
    > 1
  \]
  A série diverge
\end{frame}

%------------------------------------------------------------------------------%
\section{Lista Mínima}

%------------------------------------------------------------------------------%
\begin{frame}{Lista Mínima}
  Estudar as Seção 6.7 da Apostila

  \vspace{\baselineskip}
  Exercícios: 1a-f

  \vfill
  Atenção: A prova é baseada no livro, não nas apresentações
\end{frame}

%------------------------------------------------------------------------------%
\end{document}
%------------------------------------------------------------------------------%
